% Part: incompleteness
% Chapter: introduction
% Section: decidability

\documentclass[../../../include/open-logic-section]{subfiles}

\begin{document}

\olfileid[hi]{inc}{int}{dec}

\olsection{अनिर्णेयता और अपूर्णता}

ग्योडेल (G\"odel) के अपूर्णता प्रमेयों के प्रमाण में वाक्यविन्यास के
अंकगणितीकरण की आवश्यकता होती है। किंतु उसके बिना भी हम केवल इस मान्यता से
कुछ सुंदर परिणाम प्राप्त कर सकते हैं कि कोई सिद्धांत सभी !!{decidable}
संबंधों को !!{represents} करता है। इसका प्रमाण विराम समस्या की अनिर्णेयता के
प्रमाण जैसा एक विकर्ण तर्क है।

\begin{thm}
यदि $\Gamma$ ऐसा संगत सिद्धांत है जो प्रत्येक !!{decidable} संबंध को
!!{represents} करता है, तो $\Gamma$ !!{decidable} नहीं है।
\end{thm}

\begin{proof}
मान लें कि $\Gamma$ !!{decidable} है। हम दिखाते हैं कि यदि $\Gamma$ प्रत्येक
!!{decidable} संबंध को !!{represents} करता है, तो वह असंगत होना चाहिए।

!!^{decidable} गुणों (एकस्थानी संबंधों) का निरूपण एक मुक्त चर वाले
!!{formula} से होता है। मान लें कि $!A_0(x)$, $!A_1(x)$, \dots, ऐसे सभी
!!{formula} का संगणनीय परिगणन है। अब निम्न समुच्चय $D \subseteq \Nat$ पर
विचार करें:
\[
D = \Setabs{n}{\Gamma \Proves \lnot !A_n(\num{n})}
\]
समुच्चय $D$ !!{decidable} है, क्योंकि $n \in D$ है या नहीं, इसे जाँचने के
लिए हम पहले $!A_n(x)$ और उससे $\lnot !A_n(\num{n})$ की गणना कर सकते हैं।
स्पष्टतः $!A_n(x)$ में $x$ की प्रत्येक मुक्त उपस्थिति के स्थान पर पद
$\num{n}$ रखना और $!A(\num{n})$ के आगे $\lnot$ लगाना यांत्रिक कार्य है।
मान्यता के अनुसार $\Gamma$ !!{decidable} है, अतः हम जाँच सकते हैं कि $\lnot
!A(\num{n}) \in \Gamma$ है या नहीं। यदि है, तो $n \in D$; यदि नहीं, तो $n
\notin D$। इसलिए $D$ भी !!{decidable} है।

चूँकि $\Gamma$ सभी !!{decidable} गुणों को !!{represents} करता है, वह~$D$ को
भी !!{represents} करता है। और~$\Gamma$ में $D$ को !!{represents} करने वाले
!!{formula}, $!A_0(x)$, $!A_1(x)$, \dots, में सम्मिलित हैं। अतः कोई संख्या
$d$ ऐसी मानें कि $!A_d(x)$,~$\Gamma$ में $D$ को !!{represents} करता है। यदि
$d \notin D$, तो चूँकि $!A_d(x)$,~$D$ को !!{represents} करता है, $\Gamma
\Proves \lnot !A_d(\num{d})$। किंतु इसका अर्थ है कि $d$,~$D$ की परिभाषक
शर्त पूरी करता है, और इसलिए $d \in D$। यह $d \notin D$ का विरोधाभास है। अतः
परोक्ष प्रमाण से $d \in D$।

चूँकि $d \in D$,~$D$ की परिभाषा के अनुसार $\Gamma \Proves \lnot
!A_d(\num{d})$। दूसरी ओर, चूँकि $!A_d(x)$, $\Gamma$ में~$D$ को
!!{represents} करता है, $\Gamma \Proves !A_d(\num{d})$। अतः $\Gamma$ असंगत
है।
\end{proof}

\begin{explain}
पिछला प्रमेय दिखाता है कि सभी !!{decidable} संबंधों को !!{represents} करने
वाला कोई संगत सिद्धांत !!{decidable} नहीं हो सकता। हम दिखाएँगे कि $\Th{Q}$
सभी !!{decidable} संबंधों को !!{represents} करता है; इसका अर्थ है कि $\Th{Q}$
को समाहित करने वाले सभी सिद्धांत, जैसे $\Th{PA}$ और $\Th{TA}$, भी ऐसा करते
हैं और इसलिए वे भी !!{decidable} नहीं हैं। (चूँकि ये सभी सिद्धांत मानक
प्रतिरूप में सत्य हैं, वे सभी संगत हैं।)

इस परिणाम से हम प्रथम अपूर्णता प्रमेय का एक दुर्बल रूप भी प्राप्त कर सकते
हैं। प्रत्येक !!{axiomatizable} और !!{complete} सिद्धांत !!{decidable} है।
अतः ऐसे संगत सिद्धांत जो !!{axiomatizable} हैं और सभी !!{decidable} गुणों को
!!{represents} करते हैं, !!{complete} नहीं हो सकते।
\end{explain}

\begin{thm}
यदि $\Gamma$ !!{axiomatizable} और !!{complete} है, तो वह !!{decidable} है।
\end{thm}

\begin{proof}
प्रत्येक असंगत सिद्धांत !!{decidable} है, क्योंकि असंगत सिद्धांत सभी
!!{sentence} समाहित करते हैं; इसलिए ``क्या $!A \in \Gamma$ है?'' प्रश्न का
उत्तर सदा ``हाँ'' होता है, अर्थात् उसका निर्णय किया जा सकता है।

अब मान लें कि $\Gamma$ संगत है और साथ ही !!{axiomatizable} तथा !!{complete}
है। चूँकि $\Gamma$ !!{axiomatizable} है, वह !!{computably enumerable} है।
वास्तव में हम किसी संगणनीय फलन से~$\Gamma$ के अभिगृहीतों की सभी सही
!!{derivation} का परिगणन कर सकते हैं। किसी सही !!{derivation} से हम उसके द्वारा
!!{derive} किया गया !!{sentence} गणित कर सकते हैं; अतः दोनों को मिलाकर ऐसा
संगणनीय फलन मिलता है जो~$\Gamma$ के सभी प्रमेयों का परिगणन करता है। चूँकि
$\Gamma$ संगत और !!{complete} है, कोई !!{sentence},~$\Gamma$ का प्रमेय है यदि
और केवल यदि $\lnot !A$ प्रमेय नहीं है। अतः हम निम्न प्रकार तय कर सकते हैं कि
$!A \in \Gamma$ है या नहीं। $\Gamma$ के सभी प्रमेयों का परिगणन करें। जब इस
सूची में $!A$ आए, तब हमें ज्ञात होगा कि $\Gamma \Proves !A$। जब सूची में
$\lnot !A$ आए, तब हमें ज्ञात होगा कि $\Gamma \Proves/ !A$। चूँकि $\Gamma$
!!{complete} है, इनमें से एक स्थिति अंततः अवश्य आएगी, इसलिए प्रक्रिया अंततः
उत्तर देगी।
\end{proof}

\begin{cor}
\ollabel{cor:incompleteness}
यदि $\Gamma$ संगत और !!{axiomatizable} है तथा प्रत्येक !!{decidable} गुण को
!!{represents} करता है, तो वह !!{complete} नहीं है।
\end{cor}

\begin{proof}
यदि $\Gamma$ !!{complete} होता, तो पिछले प्रमेय के अनुसार वह !!{decidable}
होता (क्योंकि वह !!{axiomatizable} और संगत है)। किंतु चूँकि $\Gamma$ प्रत्येक
!!{decidable} गुण को !!{represents} करता है, पहले प्रमेय के अनुसार वह
!!{decidable} नहीं है।
\end{proof}

\begin{prob}
दिखाएँ कि $\Th{TA} = \Setabs{!A}{\Sat{N}{!A}}$ !!{axiomatizable} नहीं है।
आप मान सकते हैं कि $\Th{TA}$ सभी निर्णेय गुणों को निरूपित करता है।
\end{prob}

जब हम स्थापित कर लें कि, उदाहरणतः, $\Th{Q}$ सभी !!{decidable} गुणों को
!!{represents} करता है, तब उपप्रमेय बताता है कि $\Th{Q}$ अपूर्ण होना चाहिए।
किंतु इसका प्रमाण किसी स्वतंत्र !!{sentence} का उदाहरण नहीं देता; वह केवल
दिखाता है कि ऐसा !!a{sentence} अवश्य विद्यमान है। उदाहरण पाने के लिए हमें
वाक्यविन्यास का अंकगणितीकरण करके ग्योडेल (G\"odel) के मूल प्रमाण-विचार का
अनुसरण करना होगा। और निस्संदेह, हमें अब भी पहला दावा सिद्ध करना है कि वास्तव
में $\Th{Q}$ सभी !!{decidable} गुणों को !!{represents} करता है।

ध्यान रहे कि प्रत्येक \emph{रोचक} सिद्धांत अपूर्ण या अनिर्णेय नहीं होता। ऐसे
अनेक सिद्धांत हैं जो रोचक गणितीय तथ्यों का वर्णन करने के लिए पर्याप्त प्रबल
हैं, पर ग्योडेल (G\"odel) के परिणाम की शर्तें पूरी नहीं करते। उदाहरणतः,
$\Th{Pres} = \Setabs{!A \in \Lang{L_{A^+}}}{\Sat{N}{!A}}$, अर्थात् मानक
प्रतिरूप में सत्य,~$\times$ रहित अंकगणितीय भाषा के !!{sentence} का समुच्चय,
पूर्ण और निर्णेय दोनों है। यह सिद्धांत प्रेस्बर्गर अंकगणित कहलाता है और
प्राकृतिक संख्याओं के वे सभी सत्य सिद्ध करता है जिन्हें केवल $\Obj{0}$,
$\prime$ और~$+$ से सूत्रित किया जा सकता है।

\end{document}
